\documentclass[../NDCNNAI_main.tex]{subfiles}
\graphicspath{{\subfix{../fig/}}}

\begin{document}
\subsection{Постановка задачи обучения FMLP}
Введём модель функционального многослойного перцептрона (FMLP). 
Пусть каждое наблюдение задаётся парой $(G, T)$, где $G$ --- функциональный предиктор (принадлежащий некоторому функциональному пространству, например, $C(Z,\mathbb{R})$ с компактным $Z \subset \mathbb{R}^u$), а $T \in \mathbb{R}^o$ --- целевая переменная. 
Рассматривается параметрическая модель вида
\[
H(w, g) = U\!\left(w_0, \int_Z F_1(w_1, x)g(x)d\mu(x), \dots, \int_Z F_k(w_k, x)g(x)d\mu(x)\right),
\]
где $w = (w_0, w_1, \dots, w_k)$ --- параметры, $F_\ell(w_\ell, \cdot)$ --- параметрические ядра (например, малые нейронные сети или базисные разложения), $U$ --- внешняя ограниченная равномерно непрерывная функция (например, обычный MLP), $\mu$ --- конечная борелевская мера.

В процессе обучения мы располагаем выборкой $\{(G_i, T_i)\}_{i=1}^n$, причём каждая функция $G_i$ известна лишь по дискретным измерениям в точках $X_{ij}$:
\[
Y_{ij} = G_i(X_{ij}) + \varepsilon_{ij}, \quad j=1,\dots,m_i.
\]
Требуется оценить параметры $w$, минимизирующие ожидаемый риск
\[
\lambda(w) = \mathbb{E}\, c\bigl(T, H(w, G)\bigr),
\]
где $c$ --- непрерывная функция потерь. 
На практике мы минимизируем эмпирический риск, в котором интегралы заменены на выборочные средние:
\[
\hat{\lambda}_{n,m}(w) = \frac{1}{n}\sum_{i=1}^n c\!\left(T_i, U\!\left(w_0, \frac{1}{m_i}\sum_{j=1}^{m_i} F_1(w_1, X_{ij})Y_{ij}, \dots, \frac{1}{m_i}\sum_{j=1}^{m_i} F_k(w_k, X_{ij})Y_{ij}\right)\right).
\]
Обозначим через $\hat{w}_{n,m}$ любой минимизатор $\hat{\lambda}_{n,m}(w)$ на компактном множестве параметров $W$.

\begin{theorem}[О согласованности эмпирического риск минимизатора (ERM) для FMLP]
Предположим, что выполнены следующие условия:
\begin{enumerate}
    \item[(A1)] Параметрические ядра $F_\ell(w_\ell, \cdot)$ измеримы по $x$, непрерывны по $w_\ell$ и равномерно ограничены интегрируемой функцией: $|F_\ell(w_\ell, x)| \le d_\ell(x)$, $d_\ell \in L^p(\mu)$.
    \item[(A2)] Входные функции $G_i$ принадлежат $L^q(\mu)$ почти наверное, $\mathbb{E}\|G\|_q < \infty$ ($q$ сопряжено к $p$).
    \item[(A3)] Внешняя функция $U$ ограничена и равномерно непрерывна на соответствующем множестве.
    \item[(A4)] Функция потерь $c$ непрерывна и удовлетворяет условию мажорирования: $|c(t,y)| \le h(t) + \gamma|y|$, $\mathbb{E}h(T) < \infty$.
    \item[(A5)] Точки измерений $X_{ij}$ независимы, одинаково распределены с законом $\mu$, ошибки $\varepsilon_{ij}$ центрированы, независимы и имеют конечный первый момент.
\end{enumerate}
Тогда эмпирические минимизаторы $\hat{w}_{n,m}$:
\[
\lim_{n\to\infty} \lim_{m\to\infty} \operatorname{dist}\!\bigl(\hat{w}_{n,m}, W^*\bigr) = 0 \quad \text{п.н.},
\]
где $W^* = \operatorname{arg\,min}_{w\in W} \lambda(w)$ --- множество истинных минимизаторов эмпирического риска, а $\operatorname{dist}$ --- расстояние Хаусдорфа.
\end{theorem}

\begin{proof}[Эскиз доказательства]
Доказательство разбивается на несколько шагов.

\textbf{1. Равномерная сходимость монте-карловских приближений интегралов.}
Для каждого фиксированного $i$ и $\ell$ рассмотрим
\[
M_\ell^{(i)}(w_\ell) = \frac{1}{m_i}\sum_{j=1}^{m_i} F_\ell(w_\ell, X_{ij})Y_{ij}, \quad I_\ell^{(i)}(w_\ell) = \int F_\ell(w_\ell, x)G_i(x)d\mu(x).
\]
Используя предположения (A1), (A2), (A5) и теорему о равномерном усиленном законе больших чисел, получаем
\[
\sup_{w_\ell \in W_\ell} \bigl|M_\ell^{(i)}(w_\ell) - I_\ell^{(i)}(w_\ell)\bigr| \xrightarrow[m_i\to\infty]{\text{п.н.}} 0.
\]
Это следует из того, что класс функций $\{x \mapsto F_\ell(w_\ell, x)G_i(x): w_\ell \in W_\ell\}$ является гладким классом, а шумовая компонента также удовлетворяет УЗБЧ.

\textbf{2. Равномерная сходимость эмпирического риска к популяционному.}
Определим <<идеальный>> эмпирический риск
\[
\bar{\lambda}_n(w) = \frac{1}{n}\sum_{i=1}^n c\bigl(T_i, H(w, G_i)\bigr).
\]
Из равномерной непрерывности $U$ и $c$ и предыдущего шага выводим
\[
\sup_{w\in W} \bigl|\hat{\lambda}_{n,m}(w) - \bar{\lambda}_n(w)\bigr| \xrightarrow[n\to\infty]{\text{п.н.}} 0.
\]
Далее, класс функций $\{(g,t) \mapsto c(t, H(w, g)) : w \in W\}$ также является гладким классом (благодаря компактности $W$, доминированию и непрерывности), поэтому
\[
\sup_{w\in W} \bigl|\bar{\lambda}_n(w) - \lambda(w)\bigr| \xrightarrow[n\to\infty]{\text{п.н.}} 0.
\]
Комбинируя, получаем
\[
\sup_{w\in W} \bigl|\hat{\lambda}_{n,m}(w) - \lambda(w)\bigr| \xrightarrow[n\to\infty,\; m\to\infty]{\text{п.н.}} 0.
\]

\textbf{3. Сходимость минимизаторов.}
В силу компактности $W$ и непрерывности $\lambda(w)$ множество $W^*$ непусто. 
Из равномерной сходимости функций и леммы о непрерывности argmin-отображения следует
\[
\operatorname{dist}\!\bigl(\operatorname{arg\,min} \hat{\lambda}_{n,m}, W^*\bigr) \xrightarrow[n\to\infty,\; m\to\infty]{\text{п.н.}} 0.
\]
\end{proof}

\subsection{Практическая реализация: дискретизация, выбор базиса, регуляризация}
На практике функциональные данные заданы дискретно, поэтому необходимо аппроксимировать интегралы. 
Основные подходы:
\begin{itemize}
    \item \textbf{Дискретизация по сетке.} Если функция $g$ задана на равномерной сетке $t_1,\dots,t_m$, то интеграл заменяется суммой
    \[
    \int_Z F(w, x)g(x)d\mu(x) \approx \sum_{j=1}^m \omega_j F(w, t_j)g(t_j),
    \]
    где $\omega_j$ --- веса квадратурной формулы (например, прямоугольников, трапеций).
    \item \textbf{Метод Монте-Карло.} Если точки $X_j$ выбраны случайно согласно мере $\mu$, то используется оценка
    \[
    \frac{1}{m}\sum_{j=1}^m F(w, X_j)g(X_j).
    \]
    Это особенно удобно при нерегулярных измерениях.
\end{itemize}

\textbf{Выбор базиса для ядер $F_\ell$.} 
Чтобы обеспечить универсальную аппроксимацию, необходимо, чтобы семейство $\{F_\ell(w_\ell, \cdot)\}$ было плотно в $L^q(\mu)$. 
Распространённые варианты:
\begin{itemize}
    \item \textbf{B-сплайны.} Локальность и гладкость, легко налагать ограничения (например, носитель в $[-T,0]$ для причинности).
    \item \textbf{Ортогональные полиномы (Лежандра, Чебышева).} Этот вариант обеспечивает хорошую численная устойчивость, возможность точного интегрирования.
    \item \textbf{Вейвлеты.} Многомасштабный анализ, полезно для сигналов с особенностями.
    \item \textbf{Маленькие MLP.} Небольшая нейросеть, принимающая на вход $x$, обеспечивает универсальную аппроксимацию на компакте.
\end{itemize}

\textbf{Регуляризация.} 
Для улучшения обобщения и обеспечения устойчивости в динамических задачах применяют:
\begin{itemize}
    \item \textbf{Гребневую регуляризацию (weight decay)} на параметры $w$.
    \item \textbf{Ограничение нормы ядер} $\|F_\ell(w_\ell, \cdot)\|_{L^q}$, что контролирует липшицевость оператора.
    \item \textbf{Специальные ограничения для причинности:} например, требовать $F_\ell(w_\ell, \tau)=0$ при $\tau>0$, если окно $[-T,0]$ соответствует прошлому.
\end{itemize}

\subsection{Интерпретация FMLP как функциональной надстройки над моделями I–IV}
FMLP естественно встраивается в эту структуру:
\begin{itemize}
    \item \textbf{Model I (NARX).} Обычный MLP, принимающий конечный вектор регрессоров. 
    FMLP обобщает это, заменяя конечный вектор на функциональное окно $g(\tau) = (u(\tau), y(\tau))$, $\tau \in [-T,0]$. 
    Тогда предсказание имеет вид
    \[
    \hat{y}(k+1) = \Phi\bigl(z(k)\bigr), \quad z(k) = \Bigl(\int F_1(w_1,\tau)g_k(\tau)d\tau, \dots, \int F_k(w_k,\tau)g_k(\tau)d\tau\Bigr).
    \]
    \item \textbf{Model II (State-space NN).} Функциональный слой извлекает признаки $z(k)$, которые затем подаются в динамику состояния:
    \[
    x(k+1) = \Psi\bigl(z(k)\bigr), \quad \hat{y}(k) = \Theta\bigl(x(k)\bigr).
    \]
    \item \textbf{Model III (Free-run).} После идентификации модель запускается в замкнутом контуре, где вместо истинных $y$ используются предсказанные $\hat{y}$. 
    Устойчивость можно обеспечивать регуляризацией Якобиана или ISS-критериями.
    \item \textbf{Model IV (Teacher-forcing).} Идентификация в разомкнутом контуре с измеренны-
ми y.
\end{itemize}


\subsection{Применение к идентификации и управлению динамическими системами}
\begin{itemize}
    \item \textbf{Идентификация.} FMLP позволяет обучать модели напрямую по нерегулярным выборкам, без предварительного восстановления траекторий. 
    Благодаря теореме о согласованности, оценка параметров даёт состоятельную модель при увеличении объёма данных.
    \item \textbf{Управление.} В схемах предсказывающего управления модель, обученная как функциональный многослойный перцептрон, может использоваться для прогноза будущего поведения системы. 
    Функциональное представление окна прошлого позволяет учесть историю без роста размерности вектора состояния.
    \item \textbf{Устойчивость.} Применяя регуляризацию, описанную выше, можно гарантировать BIBO-устойчивость (bounded input bounded output) или ISS-свойства (input to state stability) замкнутой системы.
    \item \textbf{Инвариантность к сдвигу.} Поскольку функциональные ядра $F_\ell$ работают с окном фиксированной длины, модель автоматически инвариантна к временн`ым сдвигам в пределах этого окна.
\end{itemize}

\begin{remarks}
    \leavevmode
    \begin{itemize}
        \item Теорема о согласованности гарантирует, что при увеличении числа наблюдений и числа точек измерений на каждой функции, оценённые параметры сходятся к оптимальным. 
        Это основа для надёжного применения FMLP в практических задачах.
        \item Выбор базиса и регуляризации должен учитывать специфику задачи: для гладких сигналов подходят сплайны, для сигналов с разрывами --- вейвлеты, для задач с необходимостью интерпретации --- ортогональные полиномы.
        \item FMLP естественным образом расширяет классические модели идентификации, позволяя работать непосредственно с функциональными данными, что особенно важно при нерегулярной выборке или при наличии пропусков.
    \end{itemize}
\end{remarks}

Для углубления в тему, можно заглянуть в следующие источники \cite{Rossi2005, Stinchcombe1999, Hornik1991, Ramsay2005}.

\end{document}